\section{Experiments}
\label{sec:experiments}

\paragraph{Setup.} We evaluate Temp-ViT in two complementary settings
that the literature treats separately. \textbf{(a) Sentence-level
hallucination detection.} Given a long-form response decomposed into
claim- or sentence-level units, the task is to assign an uncertainty
score to each unit that aligns with a per-unit hallucination label
obtained from atomic-fact verification. Table~\ref{tab:main} reports
results on the three corpora that anchor this task in recent work --
RAGTruth \citep{niu2024ragtruth}, FactScore \citep{min2023}, and
LongFact \citep{wei2024}. \textbf{(b) Sequence-level uncertainty on
LM-Polygraph.} Given a whole response, the task is to assign a single
uncertainty score whose ranking of responses aligns with response-level
quality. Table~\ref{tab:polygraph} reports results on the LM-Polygraph
benchmark \citep{fadeeva2023lmpolygraph,vashurin2025lmpolygraph},
covering question answering (QA: TriviaQA, CoQA, MMLU, GSM8k),
neural machine translation (NMT: WMT14 Fr--En, WMT19 De--En), and
summarization (SUM: XSUM). Both setups are run on the LM-Polygraph
infrastructure so that the baselines below share an identical
generation, sampling, and scoring pipeline.

We attempt to answer:
(Q\textsubscript{1}) Does Temp-ViT outperform sequence-level and
consistency-based baselines for sentence-level hallucination
detection? (Q\textsubscript{2}) Does its per-token output remain
competitive when aggregated to a single sequence-level score on the
heterogeneous Polygraph tasks? (Q\textsubscript{3}) How do the gains
scale with model capacity, from 7B to 70B?

\paragraph{LLMs.} All experiments are run on three instruction-tuned,
open-weight LLMs:
Mistral-7B-Instruct-v0.2 \citep{albertqjiang2023} (Mis-7B-Inst),
Llama-3.1-8B-Instruct \citep{touvron2023} (LlaMa-8B-Inst), and
Llama-3.1-70B-Instruct \citep{touvron2023} (LlaMa-70B-Inst). For each
LLM we collect post-residual hidden states across all transformer
layers and all generated tokens, yielding the activation tensor
$A \in \mathbb{R}^{L \times T \times D}$ that drives our method
(Section~\ref{sec:notation}).

\paragraph{Datasets.} For sentence-level hallucination detection
(Table~\ref{tab:main}) we use the standard splits of
\textbf{RAGTruth} \citep{niu2024ragtruth} -- a retrieval-augmented
generation corpus with span-level hallucination annotations across
QA, summarization, and data-to-text prompts;
\textbf{FactScore} \citep{min2023} -- atomic-fact-verified biography
generations; and \textbf{LongFact} \citep{wei2024} -- long-form
factuality prompts paired with the SAFE atomic-fact verifier of the
same paper. For each corpus we follow the original protocol to obtain
per-claim or per-sentence support labels, and use them as the
auxiliary supervision signal $\tilde{y}_m$ introduced in
Section~\ref{sec:notation}. For LM-Polygraph (Table~\ref{tab:polygraph})
we follow \citet{vashurin2025lmpolygraph} for subset selection, prompt
formatting, and few-shot example sourcing; quality measures for PRR
are accuracy for short-form QA, AlignScore for long-form QA and
summarization, and COMET for translation, exactly as in the reference.

\paragraph{Baselines.} We compare against the set of baselines
maintained by LM-Polygraph
\citep{fadeeva2023lmpolygraph,vashurin2025lmpolygraph}, covering
\emph{information-based} (Sequence Probability~SP, Perplexity~PPL,
Mean Token Entropy~MTE, Monte Carlo Sequence Entropy~MCSE, and
Monte Carlo Normalized Sequence Entropy~MCNSE),
\emph{consistency-based} (Degree Matrix~DegMat, Sum of
Eigenvalues of the Graph Laplacian~EigValLaplacian),
\emph{hybrid} (Semantic Entropy, SAR), \emph{verbalized}
(P(True)), and the consistency-with-token-confidence CoCoA family.
For sentence-level hallucination, each sequence-level baseline is
re-applied per claim or per sentence; for LM-Polygraph each baseline
is used as published. Our own rows split between Stage-2 Temp-ViT
(per-window output aggregated by top-$k$ MIL into a sequence-level
score, Section~\ref{sec:temporal_head}) and Stage-3 Temp-ViT
(per-atom output trained against $\tilde{y}_m$,
Section~\ref{sec:stage3}); Stage 2 is the natural comparison for
Table~\ref{tab:polygraph}, while Stage 3 is what shows up in
Table~\ref{tab:main}.

\paragraph{Evaluation Measure.} As the primary metric in both tables
we report the Prediction Rejection Ratio
\citep{malinin2021uncertainty},
\begin{equation}
    \mathrm{PRR} = \frac{\mathrm{AUC}_{\mathrm{unc}} - \mathrm{AUC}_{\mathrm{rnd}}}{\mathrm{AUC}_{\mathrm{oracle}} - \mathrm{AUC}_{\mathrm{rnd}}},
    \label{eq:prr}
\end{equation}
which measures how well an uncertainty score ranks predictions by
quality relative to an oracle that perfectly ranks them. We compute
PRR up to a rejection threshold of 50\% to prevent degenerate inflation
at near-empty retained sets. AUROC, ECE, and Brier scores under the
same setup are reported in Appendix~\ref{app:extended} as an extended
view of Table~\ref{tab:main}.

\paragraph{Generation setup.} For both tables we report results under
greedy decoding. Stochastic-sampling and best-sample variants follow
the LM-Polygraph defaults and are reported in
Appendix~\ref{app:extended}.

Next we present the main results, with extended experimental details
in Appendix~\ref{app:extended}.
